\section{习题课 \#2（9月20日）}

\subsection{余子式、伴随矩阵与 Cramer 法则}

\begin{definition}
设 $A=(a_{ij})\in\F^{n\times n}$。删去第 $i$ 行和第 $j$ 列所得子式记为
\begin{equation}
 M_{ij},
\end{equation}
相应的代数余子式为
\begin{equation}
 C_{ij}=(-1)^{i+j}M_{ij}.
\end{equation}
由 $C_{ij}$ 组成的矩阵 $C=(C_{ij})$ 的转置称为 $A$ 的伴随矩阵，记为
\begin{equation}
 \adj(A)=C^T.
\end{equation}
\end{definition}

\begin{theorem}
对任意 $1\leq i,k\leq n$，有
\begin{equation}
 \sum_{j=1}^n a_{kj}C_{ij}=\delta_{ik}\det A.
\end{equation}
因此
\begin{equation}
 A\adj(A)=\adj(A)A=(\det A)I_n.
\end{equation}
若 $\det A\neq0$，则
\begin{equation}
 A^{-1}=\frac{1}{\det A}\adj(A).
\end{equation}
\end{theorem}

\begin{corollary}[Cramer 法则]
若 $A\in\F^{n\times n}$ 且 $\det A\neq0$，则方程组 $Ax=b$ 有唯一解，并且
\begin{equation}
 x_j=\frac{\det A_j(b)}{\det A},
 \qquad j=1,\ldots,n,
\end{equation}
其中 $A_j(b)$ 表示将 $A$ 的第 $j$ 列替换为 $b$ 后所得矩阵。
\end{corollary}

\subsection{两类 Toeplitz 行列式}

\begin{exercise}[三对角 Toeplitz 行列式]
计算
\begin{equation}
 D_n=\begin{vmatrix}
 a&b&&&\\
 c&a&b&&\\
 &c&a&\ddots&\\
 &&\ddots&\ddots&b\\
 &&&c&a
 \end{vmatrix}_{n\times n}.
\end{equation}
\end{exercise}

\begin{solution}
沿第一行展开得到递推式
\begin{equation}
 D_n=aD_{n-1}-bcD_{n-2},
 \qquad D_0=1,
 \qquad D_1=a.
\end{equation}
设 $\eta_1,\eta_2$ 为
\begin{equation}
 t^2-at+bc=0
\end{equation}
的两个根。若 $\eta_1\neq\eta_2$，则
\begin{equation}
 D_n=\frac{\eta_1^{n+1}-\eta_2^{n+1}}{\eta_1-\eta_2}.
\end{equation}
若 $\eta_1=\eta_2=\eta$，则取极限得
\begin{equation}
 D_n=(n+1)\eta^n.
\end{equation}
\end{solution}

\begin{exercise}[常上、下三角 Toeplitz 行列式]
计算
\begin{equation}
 D_n=\det\begin{pmatrix}
 a&b&b&\cdots&b\\
 c&a&b&\cdots&b\\
 c&c&a&\cdots&b\\
 \vdots&\vdots&\vdots&\ddots&\vdots\\
 c&c&c&\cdots&a
 \end{pmatrix}.
\end{equation}
\end{exercise}

\begin{solution}
将相邻两行作差，或分别沿首行、末行处理，可得到
\begin{equation}
 D_n=(a-b)D_{n-1}+b(a-c)^{n-1}
\end{equation}
以及对称的递推式
\begin{equation}
 D_n=(a-c)D_{n-1}+c(a-b)^{n-1}.
\end{equation}
当 $b\neq c$ 时，解为
\begin{equation}
 D_n=\frac{b(a-c)^n-c(a-b)^n}{b-c}.
\end{equation}
当 $b=c$ 时，由极限或直接递推可得
\begin{equation}
 D_n=(a-b)^{n-1}\bigl(a+(n-1)b\bigr).
\end{equation}
\end{solution}

\subsection{Cauchy 行列式与循环矩阵}

\begin{theorem}[Cauchy 行列式]
设 $x_i+y_j\neq0$，且 $x_i$ 两两不同、$y_j$ 两两不同，则
\begin{equation}
 \det\left(\frac{1}{x_i+y_j}\right)_{i,j=1}^n
 =
 \frac{
 \displaystyle\prod_{1\leq i<j\leq n}(x_j-x_i)(y_j-y_i)}
 {\displaystyle\prod_{i=1}^n\prod_{j=1}^n(x_i+y_j)}.
\end{equation}
\end{theorem}

\begin{proof}
将第 $j$ 列与第一列作适当差分，可逐步提出因子 $y_j-y_1$；再对行作同样处理，可提出 $x_j-x_1$。余下行列式降为 $n-1$ 阶 Cauchy 行列式，归纳即得结论。另一种写法是将分母统一后观察分子关于 $x_i,y_j$ 分别为交错多项式，因此必含两个 Vandermonde 因子，再比较最高次项常数。
\end{proof}

\begin{theorem}[循环矩阵的特征值]
令
\begin{equation}
 C=\begin{pmatrix}
 a_0&a_1&\cdots&a_{n-1}\\
 a_{n-1}&a_0&\cdots&a_{n-2}\\
 \vdots&\vdots&\ddots&\vdots\\
 a_1&a_2&\cdots&a_0
 \end{pmatrix},
 \qquad
 f(z)=\sum_{j=0}^{n-1}a_jz^j,
\end{equation}
并取 $\omega=\e^{2\pi\ii/n}$。则
\begin{equation}
 v_k=(1,\omega^k,\omega^{2k},\ldots,\omega^{(n-1)k})^T
\end{equation}
是特征向量，对应特征值 $f(\omega^k)$。因此
\begin{equation}
 \det C=\prod_{k=0}^{n-1}f(\omega^k).
\end{equation}
\end{theorem}

\begin{exercise}[正弦行列式]
计算
\begin{equation}
 D=\det\bigl(\sin(j\theta_i)\bigr)_{i,j=1}^n.
\end{equation}
\end{exercise}

\begin{solution}
利用 Chebyshev 多项式第二类
\begin{equation}
 \sin(j\theta)=\sin\theta\,U_{j-1}(\cos\theta),
\end{equation}
且 $U_{j-1}$ 的首项系数为 $2^{j-1}$，于是
\begin{equation}
 D
 =2^{n(n-1)/2}
 \left(\prod_{i=1}^n\sin\theta_i\right)
 \left(\prod_{1\leq i<j\leq n}(\cos\theta_j-\cos\theta_i)\right).
\end{equation}
这与原稿中利用复指数和 Vandermonde 行列式的推导等价。
\end{solution}

\subsection{行列式的求导与几个加分题}

\begin{theorem}[行列式求导公式]
设 $A(t)=(a_{ij}(t))$ 的各元素可微。令 $B_j(t)$ 表示将 $A(t)$ 的第 $j$ 列替换为其导数列后所得矩阵，则
\begin{equation}
 \frac{\dd}{\dd t}\det A(t)
 =\sum_{j=1}^n\det B_j(t)
 =\sum_{i,j=1}^n a'_{ij}(t)C_{ij}(t).
\end{equation}
若 $A(t)$ 可逆，还可写为 Jacobi 公式
\begin{equation}
 \frac{\dd}{\dd t}\det A(t)
 =\det A(t)\,\tr\bigl(A(t)^{-1}A'(t)\bigr).
\end{equation}
\end{theorem}

\begin{proof}
行列式对每一列均为线性函数。对乘积形式逐列求导，恰好得到把某一列替换为导数列的 $n$ 项之和。再对每个 $\det B_j$ 按第 $j$ 列展开，即得代数余子式表达式。
\end{proof}

\begin{exercise}[带参数的 Vandermonde 行列式]
令
\begin{equation}
 D_1(y)=
 \det\begin{pmatrix}
 1&1&\cdots&1\\
 y&x_1&\cdots&x_n\\
 y^2&x_1^2&\cdots&x_n^2\\
 \vdots&\vdots&&\vdots\\
 y^n&x_1^n&\cdots&x_n^n
 \end{pmatrix}.
\end{equation}
求其关于 $y$ 的展开。
\end{exercise}

\begin{solution}
Vandermonde 公式给出
\begin{equation}
 D_1(y)=
 \prod_{j=1}^n(x_j-y)
 \prod_{1\leq i<j\leq n}(x_j-x_i).
\end{equation}
若 $e_r(x_1,\ldots,x_n)$ 表示第 $r$ 个初等对称多项式，则
\begin{equation}
 \prod_{j=1}^n(x_j-y)
 =\sum_{k=0}^n(-1)^k e_{n-k}(x_1,\ldots,x_n)y^k.
\end{equation}
因此 $y^k$ 的系数为
\begin{equation}
 (-1)^k e_{n-k}(x_1,\ldots,x_n)
 \prod_{i<j}(x_j-x_i).
\end{equation}
\end{solution}

\begin{exercise}
设
\begin{equation}
 D_2=\det\bigl(1+x_i^j\bigr)_{i,j=1}^n.
\end{equation}
证明
\begin{equation}
 D_2=
 \left(2\prod_{k=1}^n x_k-\prod_{k=1}^n(x_k-1)\right)
 \prod_{1\leq i<j\leq n}(x_j-x_i).
\end{equation}
\end{exercise}

\begin{proof}
在行列式左侧添一行一列，将其写成
\begin{equation}
 D_2=
 \det\begin{pmatrix}
 1&1&1&\cdots&1\\
 -1&x_1&x_1^2&\cdots&x_1^n\\
 \vdots&\vdots&\vdots&&\vdots\\
 -1&x_n&x_n^2&\cdots&x_n^n
 \end{pmatrix}.
\end{equation}
把第一列拆成 $(-1,1,\ldots,1)^T$ 与一个只在首项非零的向量之差，并分别使用 Vandermonde 公式，即得所述表达式。
\end{proof}

\begin{exercise}[全体元素同时加参数]
设 $A=(a_{ij})$，$J=\ones\ones^T$，证明
\begin{equation}
 \det(A+tJ)=\det A+t\sum_{i,j=1}^n C_{ij}(A).
\end{equation}
\end{exercise}

\begin{proof}
矩阵 $tJ$ 的各列相同。由行列式对列的多重线性，展开 $\det(A+tJ)$ 时，含有两列以上来自 $tJ$ 的项均为零，因此只剩常数项和一次项。一次项正是所有代数余子式之和。也可直接使用矩阵行列式引理。
\end{proof}

\begin{exercise}[参数特征方程]
求齐次方程组
\begin{equation}
\begin{cases}
 (\lambda-2)x_1-3x_2-2x_3=0,\\
 -x_1+(\lambda-8)x_2-2x_3=0,\\
 2x_1+14x_2+(\lambda+3)x_3=0
\end{cases}
\end{equation}
存在非零解时的 $\lambda$。
\end{exercise}

\begin{solution}
系数行列式为
\begin{equation}
 \det\begin{pmatrix}
 \lambda-2&-3&-2\\
 -1&\lambda-8&-2\\
 2&14&\lambda+3
 \end{pmatrix}
 = (\lambda-1)(\lambda-3)^2.
\end{equation}
故且仅当
\begin{equation}
 \lambda=1\quad\text{或}\quad\lambda=3
\end{equation}
时存在非零解。
\end{solution}

\begin{remark}
一般的参数齐次方程组
\begin{equation}
 A x=\lambda x
\end{equation}
存在非零解的参数由特征多项式 $\det(\lambda I-A)$ 的零点给出。若该多项式不恒为零，则不同参数值至多有 $n$ 个；在代数闭域中按重数恰有 $n$ 个。
\end{remark}

\subsection{分段插值与严格对角占优型行列式}

手稿最后补充了分段多项式插值。若在 $n$ 个区间上分别使用次数不超过 $s$ 的多项式，则未知系数共有 $n(s+1)$ 个。插值端点条件、区间连接处的连续性条件以及边界条件必须在数量与独立性上共同保证唯一性。分段线性插值对应 $s=1$ 且只要求函数连续；三次样条对应 $s=3$、连接处要求到二阶导数连续，还需补充两个边界条件，例如自然边界条件。

\begin{theorem}[一类 $M$-矩阵判别]
设实矩阵 $A=(a_{ij})\in\R^{n\times n}$ 满足
\begin{equation}
 a_{ii}>0,
 \qquad a_{ij}<0\quad(i\neq j),
 \qquad \sum_{j=1}^n a_{ij}>0\quad(i=1,\ldots,n).
\end{equation}
则
\begin{equation}
 \det A>0.
\end{equation}
\end{theorem}

\begin{proof}
对 $n$ 归纳。以 $a_{11}>0$ 为主元消去第一列。Schur 补的非对角元仍为负，而其每一行的元素和为
\begin{equation}
 \sum_{j=2}^n\left(a_{ij}-a_{i1}a_{11}^{-1}a_{1j}\right)
 =\sum_{j=1}^n a_{ij}
 -a_{i1}a_{11}^{-1}\sum_{j=1}^n a_{1j}>0,
\end{equation}
因为 $a_{i1}<0$ 且第一行行和为正。故 Schur 补满足同样条件，由归纳假设其行列式为正。于是
\begin{equation}
 \det A=a_{11}\det(A/a_{11})>0.
\end{equation}
\end{proof}
